\section{Challenges with AI-agent analysis}
\label{sec:challenges}

This section is maintained as a living record of unresolved or recurring issues encountered during the analysis and during analyst--agent collaboration. Because this document is an internal note rather than a final publication, we prioritize explicit tracking of weaknesses, mitigation status, and collaboration failure modes.

\subsection{Practical Codex--physicist workflow}

The working model used in this note is a constrained human-in-the-loop cycle:
\begin{enumerate}
\item the physicist defines the physics objective, nominal policy, and acceptance criteria (e.g. nominal IBU iterations, uncertainty definition, figure intent);
\item the agent implements code and note changes, executes reproducible commands, and reports exact file-path and artifact updates;
\item the physicist reviews outputs (plots, numerical consistency, physics interpretation), requests targeted corrections, and approves or rejects the iteration;
\item accepted changes are propagated to the note with explicit script provenance and residual-risk statements.
\end{enumerate}

In this model, authority is intentionally asymmetric: the physicist owns physics choices and release signoff, while the agent owns implementation throughput, consistency checks, and traceability.

\subsection{How interaction is carried out in practice}

The interaction pattern in this project is not free-form chat; it is a specification--execution loop with explicit control points:
\begin{itemize}
\item requests are usually concrete and executable (script, figure, table, section);
\item each iteration requires path-level traceability to scripts, outputs, and note sections;
\item plotting/layout feedback is treated as analysis-relevant, since figure readability affects interpretation;
\item corrections are applied incrementally and revalidated, rather than batching large unreviewed changes;
\item unresolved issues are converted into tracked follow-up tasks rather than hidden in narrative text.
\end{itemize}

This interaction discipline reduced turnaround for code and note updates, but it also exposed repeated bottlenecks in formatting consistency, evolving scope, and ambiguity in implicit expert expectations.

\subsection{Agent-perspective challenges with expert collaboration}

From the agent side, the main recurring challenges are:
\begin{itemize}
\item \textbf{Implicit physics priors}: expert expectations on acceptable binning, regularization behavior, or uncertainty presentation are often assumed rather than initially stated;
\item \textbf{Precision formatting constraints}: small figure-layout details (label offsets, legend placement, pad geometry) require many short feedback loops;
\item \textbf{Method-vs-presentation coupling}: apparent cosmetic changes can alter interpretability and therefore physics conclusions;
\item \textbf{State synchronization risk}: when scripts, outputs, and note text evolve quickly, stale references can appear unless provenance is continuously refreshed;
\item \textbf{Boundary management}: the agent can execute and refactor reliably, but interpretation authority and final scientific claims must remain analyst-owned.
\end{itemize}

These challenge classes are expected in future observable expansions and motivate stricter templating for figures, stronger configuration contracts, and checklists for note updates.

\subsection{Current limitations of the thrust-only baseline}

The present release is robust enough for a preliminary thrust result, but several limitations remain:
\begin{itemize}
\item for events with exactly 3 selected particles, the thrust algorithm selects the highest-magnitude particle as the axis rather than enumerating all possible hemisphere assignments; this approximation is justified because (i) 3-particle hadronic events are rare in $Z \to q\bar{q}$ decays (impact on thrust value typically $\lesssim 0.1\%$ for $T > 0.5$), and (ii) the response-matrix migration structure dominates uncertainty in the endpoint bins where such events concentrate;
\item hadronic-event correction uncertainty is propagated as an independent component ($\Delta_{\mathrm{had}}$, central-range mean $\approx 3.2\%$); the endpoint bins remain sensitive to correction-factor extrapolation;
\item systematic inter-bin correlations use the statistical correlation matrix as a template ($\mathrm{Cov}^{(c)}_{ij} = \Delta_c(i)\,\Delta_c(j)\,\rho^{\mathrm{stat}}_{ij}$); this is an approximation that assumes systematic and statistical correlations share the same bin-to-bin structure;
\item endpoint thrust bins remain statistically and modeling sensitive;
\item analysis scope is currently one-dimensional (thrust only), while the future target includes the thrust-family observables and potentially multidimensional unfolding.
\end{itemize}

\subsection{Persistent analysis challenges}

Several recurrent issues have been tracked through the development of this note. Plot-style consistency across scripts has been an ongoing concern, with centralized style enforcement in producer scripts now largely resolving the issue. Method-to-code traceability gaps are addressed by the workflow and provenance tables in Sec.~\ref{sec:method} and Sec.~\ref{sec:results}, with \path{workflow_summary.json} serving as the canonical numeric record. The most significant open challenge is endpoint-bin instability: the final bin ($0.99<T<1.00$) is sensitive to MC modeling, hadronic correction extrapolation, and normalization convention differences, and dominates the maximum per-bin uncertainty. This is mitigated by reporting central-window metrics ($0.60\le T<0.95$) alongside full-range figures throughout the note.

\subsection{Update protocol for this section}

For each major workflow iteration, this section should record:
\begin{itemize}
\item newly discovered persistent issues;
\item whether previous issues moved from Open to In progress/Resolved;
\item concrete next action and code location for unresolved items.
\end{itemize}

Keeping this log in the main body is intentional: it makes assumptions and remaining risks visible to internal readers and future maintainers.
